% \section{PROBLEM DEFINITION}
% \label{sec:problem_definition}

% \subsection{Formalization of Real-Time Federated Multi-Task Learning for NLU}
% We consider a RT-FMTL setting for NLU, where training proceeds over communication rounds $t=1,2,\dots$ \cite{mcmahan2017fedavg,smith2017mocha,chen2024fedbone}. At each round, a subset of clients is selected, performs local updates on private data streams, and sends model updates to a central server for aggregation. Unlike purely offline FL, the client data distribution may evolve over time, so both local objectives and optimal parameters are time-dependent.

% Let $\mathcal{C}=\{1,\dots,C\}$ be the set of clients and $\mathcal{K}=\{1,\dots,K\}$ be the set of NLU tasks. Each client $c\in\mathcal{C}$ is associated with one primary task $\tau(c)\in\mathcal{K}$ (or a small task subset), and may receive new samples continuously from a local stream. The goal is to learn a federated model that preserves cross-task transfer while maintaining task-specific quality under non-IID and time-varying data.

% \subsection{Data and Task Heterogeneity Across Clients}
% Heterogeneity appears in two forms:
% \begin{itemize}
%     \item Data heterogeneity: for clients $c_i,c_j$, local distributions differ, i.e., $P_{c_i}(x,y)\neq P_{c_j}(x,y)$ (label skew, domain shift, vocabulary differences).
%     \item Task heterogeneity: clients may optimize different losses and output spaces based on $\tau(c)$, leading to mismatched gradients across tasks.
% \end{itemize}

% In addition, client availability and update frequency are asynchronous in real-time settings, which introduces system heterogeneity and stale updates \cite{mcmahan2017fedavg,li2020fedprox,chen2024fedbone}. Therefore, the optimization process must be robust to partial participation and non-uniform local progress.

% \subsection{Optimization Objective for Federated Multi-Task Learning}
% Let $\ell_{\tau(c)}(\theta_e,\theta_{\tau(c)};\xi)$ be the task-specific loss for sample $\xi\in\mathcal{D}_c^t$. The global objective at round $t$ is:
% \begin{equation}
% \begin{aligned}
% \min_{\theta_e,\{\theta_k\}} \, F_t(\Theta)\section{PROBLEM DEFINITION}
% \label{sec:problem_definition}

% \subsection{Formalization of Real-Time Federated Multi-Task Learning for NLU}
% We consider a RT-FMTL setting for NLU, where training proceeds over communication rounds $t=1,2,\dots$ \cite{mcmahan2017fedavg,smith2017mocha,chen2024fedbone}. At each round, a subset of clients is selected, performs local updates on private data streams, and sends model updates to a central server for aggregation. Unlike purely offline FL, the client data distribution may evolve over time, so both local objectives and optimal parameters are time-dependent.

% Let $\mathcal{C}=\{1,\dots,C\}$ be the set of clients and $\mathcal{K}=\{1,\dots,K\}$ be the set of NLU tasks. Each client $c\in\mathcal{C}$ is associated with one primary task $\tau(c)\in\mathcal{K}$ (or a small task subset), and may receive new samples continuously from a local stream. The goal is to learn a federated model that preserves cross-task transfer while maintaining task-specific quality under non-IID and time-varying data.

% \subsection{Data and Task Heterogeneity Across Clients}
% Heterogeneity appears in two forms:
% \begin{itemize}
%     \item Data heterogeneity: for clients $c_i,c_j$, local distributions differ, i.e., $P_{c_i}(x,y)\neq P_{c_j}(x,y)$ (label skew, domain shift, vocabulary differences).
%     \item Task heterogeneity: clients may optimize different losses and output spaces based on $\tau(c)$, leading to mismatched gradients across tasks.
% \end{itemize}

% In addition, client availability and update frequency are asynchronous in real-time settings, which introduces system heterogeneity and stale updates \cite{mcmahan2017fedavg,li2020fedprox,chen2024fedbone}. Therefore, the optimization process must be robust to partial participation and non-uniform local progress.

% \subsection{Optimization Objective for Federated Multi-Task Learning}
% Let $\ell_{\tau(c)}(\theta_e,\theta_{\tau(c)};\xi)$ be the task-specific loss for sample $\xi\in\mathcal{D}_c^t$. The global objective at round $t$ is:
% \begin{equation}
% \begin{aligned}
% \min_{\theta_e,\{\theta_k\}} \, F_t(\Theta)
% &= \sum_{c=1}^{C} p_c^t \, \mathbb{E}_{\xi\sim\mathcal{D}_c^t}
% \big[\ell_{\tau(c)}(\theta_e,\theta_{\tau(c)};\xi)\big] \\
% &\quad + \lambda \, \Omega(\{\theta_k\}),
% \end{aligned}
% \end{equation}
% where $p_c^t$ is the aggregation weight (e.g., proportional to local sample count), $\Omega$ is a regularizer encouraging stable cross-task learning, and $\lambda\ge0$ controls the sharing-specialization trade-off.

% After local updates, the server aggregates shared encoder parameters using weighted averaging,
% \begin{equation}
% \theta_e^{t+1} = \sum_{c\in\mathcal{C}_t} \alpha_c^t \, \theta_{e,c}^{t+1},
% \end{equation}
% and updates task-specific parameters by task-wise aggregation,
% \begin{equation}
% \theta_k^{t+1} = \sum_{c\in\mathcal{C}_t: \, \tau(c)=k} \beta_c^t \, \theta_{k,c}^{t+1}.
% \end{equation}
% Here, $\mathcal{C}_t$ is the participating client subset at round $t$, and $\alpha_c^t,\beta_c^t$ are normalized aggregation coefficients \cite{mcmahan2017fedavg,smith2017mocha,zhou2025mfed}. This objective captures real-time RT-FMTL by jointly optimizing shared representations and task-specific parameters under privacy and heterogeneity constraints.

% &= \sum_{c=1}^{C} p_c^t \, \mathbb{E}_{\xi\sim\mathcal{D}_c^t}
% \big[\ell_{\tau(c)}(\theta_e,\theta_{\tau(c)};\xi)\big] \\
% &\quad + \lambda \, \Omega(\{\theta_k\}),
% \end{aligned}
% \end{equation}
% where $p_c^t$ is the aggregation weight (e.g., proportional to local sample count), $\Omega$ is a regularizer encouraging stable cross-task learning, and $\lambda\ge0$ controls the sharing-specialization trade-off.

% After local updates, the server aggregates shared encoder parameters using weighted averaging,
% \begin{equation}
% \theta_e^{t+1} = \sum_{c\in\mathcal{C}_t} \alpha_c^t \, \theta_{e,c}^{t+1},
% \end{equation}
% and updates task-specific parameters by task-wise aggregation,
% \begin{equation}
% \theta_k^{t+1} = \sum_{c\in\mathcal{C}_t: \, \tau(c)=k} \beta_c^t \, \theta_{k,c}^{t+1}.
% \end{equation}
% Here, $\mathcal{C}_t$ is the participating client subset at round $t$, and $\alpha_c^t,\beta_c^t$ are normalized aggregation coefficients \cite{mcmahan2017fedavg,smith2017mocha,zhou2025mfed}. This objective captures real-time RT-FMTL by jointly optimizing shared representations and task-specific parameters under privacy and heterogeneity constraints.

\section{PROBLEM DEFINITION}  
\label{sec:problem_definition}

This section formalizes the RT-FedMTL environment for NLU. We define the mathematical framework for distributed multi-task optimization, address the challenges of data and system heterogeneity, and establish the global objective for collaborative model training across private edge clients.

\subsection{Formalization of RT-FedMTL for NLU}
We consider a RT-FedMTL environment based on NLU in which training proceeds over communication rounds $t=1,2,\dots$ \cite{mcmahan2017fedavg,smith2017mocha,chen2024fedbone}. During each round, a subset of clients is selected, local updates are performed on private data streams, and parameters are pushed to a central server for aggregation. Unlike offline FL, client data distributions may vary over time; therefore, both local goals and optimum parameters are time-dependent. Let $\mathcal{C}=\{1,\dots,C\}$ be the set of clients and $\mathcal{K}=\{1,\dots,K\}$ be the set of NLU tasks. Each client $c\in\mathcal{C}$ is associated with one primary task $\tau(c)\in\mathcal{K}$ (or a small task subset), where new samples are continuously added from a local stream. The objective is to learn a federated model that maintains cross-task transfer quality while respecting task-level specialization under non-IID and time-varying data.

\subsection{Data and Task Heterogeneity in Clients}  
Heterogeneity comes in two forms:  
\begin{itemize}  
\item Heterogeneity of data: the local distributions in clients $c_i,c_j$ will be different (i.e. $P_{c_i}(x,y)\neq P_{c_j}(x,y)$ (label skew, domain shift, vocabulary differences).  
\item Task heterogeneity: clients may optimize different losses and output spaces according to $\tau(c)$, yielding dissimilar gradients across tasks.  
\end{itemize}  
Moreover, real-time environments have client availability and update frequency asynchronous, leading to system heterogeneity and stale updates \cite{mcmahan2017fedavg,li2020fedprox,chen2024fedbone}. Thus, the optimization process must account for partial participation and non-uniform local progress.  

\subsection{Optimization Objective for Federated Multi-Task Learning}  
Let $\ell_{\tau(c)}(\theta_e,\theta_{\tau(c)};\xi)$ be the task-specific loss cost for sample $\xi\in\mathcal{D}_c^t$. The global objective at round $t$ is:  
\begin{equation}  
\begin{aligned}  
\min_{\theta_e,\{\theta_k\}} \, F_t(\Theta)  
&= \sum_{c=1}^{C} p_c^t \, \mathbb{E}_{\xi\sim\mathcal{D}_c^t}  
\big[\ell_{\tau(c)}(\theta_e,\theta_{\tau(c)};\xi)\big] \\  
&\quad + \lambda \, \Omega(\{\theta_k\}),  
\end{aligned}  
\end{equation}  
where $p_c^t$ represents the weight of the aggregation (e.g., the weight per unit local sample), $\Omega$ is a regularizer to promote a stable cross-task learning and $\lambda\ge0$ controls the sharing-specialization trade-off. After localized updates, the server aggregates common encoder parameters by weighted averaging,  
\begin{equation}  
\theta_e^{t+1} = \sum_{c\in\mathcal{C}_t} \alpha_c^t \, \theta_{e,c}^{t+1}.  
\end{equation}  
and updates custom task-specific parameters, as a result of task-wise aggregation,  
\begin{equation}  
\theta_k^{t+1} = \sum_{c\in\mathcal{C}_t: \, \tau(c)=k} \beta_c^t \, \theta_{k,c}^{t+1}.  
\end{equation}  
In $\mathcal{C}_t$, the client subset participate at round $t$ and $\alpha_c^t,\beta_c^t$ are normalized aggregation coefficients \cite{mcmahan2017fedavg,smith2017mocha,zhou2025mfed}. With privacy and heterogeneity requirements in place, the framework captures RT-FedMTL in real time by jointly optimizing both shared representations and task-specific parameters.